\documentclass[11pt, a4paper]{article}
\usepackage[T1]{fontenc}
\usepackage[margin=1in]{geometry}
\usepackage{setspace}
\usepackage{hyperref}

\onehalfspacing

\title{\textbf{Research Statement: Human-Centred Healthcare Robotics}}
\author{}
\date{}

\begin{document}

\maketitle

\section*{1. Introduction and Background}
\begin{itemize}
    \item Overview of healthcare robotics landscape
    \item Current challenges in patient care delivery
    \item Gap between technological capabilities and human-centered design
    \item Motivation for human-centered approach
\end{itemize}

\section*{2. Research Vision and Motivation}
\begin{itemize}
    \item Definition of human-centered healthcare robotics
    \item Key principles: \textit{safety, accessibility, empathy, inclusivity}
    \item Target stakeholders: patients, healthcare professionals, caregivers
    \item Alignment with healthcare demands and societal needs
\end{itemize}

\section*{3. Possible Research Directions}

\subsection*{3.1 Patient-Robot Interaction Design}
\begin{itemize}
    \item Intuitive interfaces for diverse user populations
    \item Understanding patient preferences and comfort levels
    \item Emotional and psychological impacts of robotic assistance
\end{itemize}

\subsection*{3.2 Safety and Trust in Clinical Environments}
\begin{itemize}
    \item Safety protocols for autonomous systems
    \item Building trust between patients and robots
    \item Failure modes and recovery mechanisms
    \item Ethical considerations in healthcare automation
\end{itemize}

\subsection*{3.3 Adaptive and Personalized Care Delivery}
\begin{itemize}
    \item Machine learning for individual patient profiles
    \item Real-time adaptation to patient needs
    \item Cultural sensitivity and accessibility in robot design
    \item Rehabilitation and assistive robotics
\end{itemize}

\subsection*{3.4 Collaboration Between Robots and Healthcare Professionals}
\begin{itemize}
    \item Human-robot teamwork models
    \item Cognitive load reduction for medical staff
    \item Training and integration in clinical workflows
    \item Decision support systems with human oversight
\end{itemize}

\subsection*{3.5 Accessibility and Inclusivity}
\begin{itemize}
    \item Design for elderly and disabled populations
    \item Multilingual and multicultural considerations
    \item Cost-effective solutions for diverse healthcare settings
    \item Remote care and underserved communities
\end{itemize}

\subsection*{3.6 Long-term Impact Assessment}
\begin{itemize}
    \item Longitudinal studies on patient outcomes
    \item Economic and social sustainability
    \item Policy and regulatory frameworks
    \item Environmental impact considerations
\end{itemize}

\section*{4. Research Methodology}
\begin{itemize}
    \item Interdisciplinary approach: robotics, healthcare, design, psychology
    \item Mixed-methods research: quantitative evaluation + qualitative feedback
    \item User-centered design processes with stakeholder involvement
    \item Iterative prototyping and clinical validation
\end{itemize}

\section*{5. Expected Outcomes and Impact}
\begin{itemize}
    \item Contributions to robotic system design
    \item Improved patient experiences and clinical outcomes
    \item Guidelines for ethical healthcare robotics deployment
    \item Publications and industry partnerships
\end{itemize}

\section*{6. Conclusion}
\begin{itemize}
    \item Summary of research vision
    \item Long-term goals in advancing human-centered healthcare robotics
    \item Commitment to bridging technology and human needs
\end{itemize}

\newpage
\begin{thebibliography}{99}

\bibitem{key1} Author, A. (Year). \textit{Title of Reference}. Journal Name, Volume(Issue), pages.

\bibitem{key2} Author, B. (Year). \textit{Title of Reference}. Publisher.

\end{thebibliography}

\end{document}
